% Document Type: LaTeX
% Master File: manual.tex

%%%% \chapter{\elan: the system}

This section is devoted to the user environment of the \elan{}\
language, and is partly borrowed from~\cite{BorovanskyKKMR-WRLA98}.

%%%%%%%%%%%%%%%%%%%%%%%%%%%%%%%%%%%%%%%%%%%%%%%%%%%%%%%%%%%%
\section{Global architecture}
%%%%%%%%%%%%%%%%%%%%%%%%%%%%%%%%%%%%%%%%%%%%%%%%%%%%%%%%%%%%

The \elan\ environment consists of several components, as depicted in
Figure~\ref{ELAN-env}. The preprocessor expands a few concise
constructions allowed  in the language.
The parser checks the syntax of programs and verifies that terms are
syntactically well-formed.  The interpreter is an interactive tool
allowing the user to check that the results he expects are indeed
obtained. Initially, the preprocessor and the parser were integrated
in the interpreter and so it was not possible to call them as
stand-alone processes.  The compiler transforms specifications into
independent executable \texttt{C} code (Section~\ref{compilo}). It is
a stand-alone tool without any preprocessing and parsing
facilities. The preprocessing and parsing phases are provided by the
interpreter, which communicates the result of these processes to the
compiler via a data exchange format called \REF. Section~\ref{ref_format}
describes some tools to translate \elan\ to \REF\ and vice versa.
Especially, the translation from \elan\ to \REF\ performed by the 
\texttt{query2ref} tool relies mainly on reusing the parser initially
integrated in the interpreter.

\begin{figure*}[ht] %zoom 50
\centering
\input{ELAN-env.pstex_t}
\label{ELAN-env}
\caption{Global architecture of the \elan\ environment}
\end{figure*}


%%%%%%%%%%%%%%%%%%%%%%%%%%%%%%%%%%%%%%%%%%%%%%%%%%%%%%%%%%%%
\section{The preprocessor}
%%%%%%%%%%%%%%%%%%%%%%%%%%%%%%%%%%%%%%%%%%%%%%%%%%%%%%%%%%%%

The \elan\ syntax provides a few fancy constructions to perform
textual replacements. So the preprocessor may be used to automatically
generate parts of specifications used to analyse the rest of a
program.  It should be emphasised that there is strong interaction
between the parser, the preprocessor and the interpreter. 
See Section~\ref{preprocessor} for a description of the preprocessor
functionality. 


%%%%%%%%%%%%%%%%%%%%%%%%%%%%%%%%%%%%%%%%%%%%%%%%%%%%%%%%%%%%
\section{The parser}
%%%%%%%%%%%%%%%%%%%%%%%%%%%%%%%%%%%%%%%%%%%%%%%%%%%%%%%%%%%%

The \elan\ parser is based on Earley's one~\cite{Earley} and is using sort
informations in order to determine as soon as possible in the analysis
the right choice. This explains why in the current version, some
construction of the language
require to give the sort information together with the term.

%%%%%%%%%%%%%%%%%%%%%%%%%%%%%%%%%%%%%%%%%%%%%%%%%%%%%%%%%%%%
\section{The interpreter}
%%%%%%%%%%%%%%%%%%%%%%%%%%%%%%%%%%%%%%%%%%%%%%%%%%%%%%%%%%%%


The interpreter takes a well-formed program and a well-formed query
(both checked by the parser) and applies the rules and strategies
defined in the program to the query. In order to find which rules can
apply, the selection is guided by the top symbol of the rules: only
those rules whose left-hand side has the same top symbol as the term
to be reduced are selected. They are then tried in the order given in
the program.  Another kind of choices arbitrarily made by the
interpreter is for the strategy $\dc(S_1,...,S_n)$ that should select
randomly a non-failing strategy among $S_1,...,S_n$. In practice, the
interpreter selects the first one, so implements \dc\ and \first\ in the
same way. Note however that there exists a  version of
\elan\ which concurrently executes the $n$ strategies and selects the
first one which terminates without failure~\cite{BorovanskyCastro-WRLA98}.

Once a set of rules is selected, a many-to-one matching algorithm 
is applied. When associative and
commutative  ($AC$ for short) operators are involved, an external 
one-to-one $AC$-matching algorithm described in~\cite{Eker-AcMatch-95}
is called. 
This algorithm is not fully integrated in the interpreter, so 
data structure conversions are required and lower the efficiency of
the $AC$-matching, already  quite complex.
Once a match is found, local evaluations are performed and if all 
succeed, the result term is built, taking advantage of the right-hand 
side of the rule and of term sharing.


%%%%%%%%%%%%%%%%%%%%%%%%%%%%%%%%%%%%%%%%%%%%%%%%%%%%%%%%%%%%
\section{The compiler}
\label{compilo}
%%%%%%%%%%%%%%%%%%%%%%%%%%%%%%%%%%%%%%%%%%%%%%%%%%%%%%%%%%%%

The \elan{} compiler transforms a logic description into an executable
binary file.  The compiler (for detailed description,
see~\cite{Vittek-RTA96,MoreauK-PLILP+ALP98}) produces an efficient
\verb|C| code, which is later compiled by the \verb|cc| or the
\verb|gnu gcc| compiler. 
The executable binary code is dependent on the particular
architecture, because a small part of the \elan{} library performing
the basic non-deterministic operations is written in assembly
language~\cite{Moreau-JICSLPW98}.
Up to now, this is the only reason that makes
the compiler available only on the architectures DEC-ALPHA, SUN4 and
Intel-PC.  
%This release of the \elan{} compiler does not compile the
%whole \elan{} language: the main restriction is that the \elan{}
%compiler does not treat Input/Output and reflective capabilities of
%the interpreter. There is also a restriction on the class of rules
%containing AC-symbols that can be compiled.

\subsection{Non-linear rules}
%%%%%%%%%%%%%%%%%%%%%%%%%%%%%%%%%%%%%%%%%%%%%%%%%%%%%%%%%%%%
The \elan\ compiler does not support non-left-linear rules (a rule is
called non-left-linear if a variable occurs more than once in the
left-hand side).
The many-to-one matching algorithm only compiles linear rules. 
The \elan\ parser automatically transforms non-left-linear rules into
left-linear conditional rules by renaming variables that occurs more
than once and adding some conditions to ensure their equality.

For example, the rule \texttt{f(x,x) => g(x)} is transformed into
\texttt{f(x,y) => g(x) if x==y}.


\subsection{Built-ins}
%%%%%%%%%%%%%%%%%%%%%%%%%%%%%%%%%%%%%%%%%%%%%%%%%%%%%%%%%%%%
To be more efficient, builtin data type such as
\textit{builtinInt} and \textit{identifier} are compiled in a
particular way. 

The main drawback of this approach is that only completely defined
functions can be defined on builtin data type.

Let \texttt{f(@) : (builtinInt) builtinInt} be an operator. 
For each integer $n$, $f(n)$ must be reducible to an integer. Suppose
defined the rule 
\texttt{[] f(n) => n-1 if n > 1}, the term \texttt{f(0)} can not be
reduced to a builtin integer. In this case, behaviour is undefined.
This explains why \textit{builtinInt} are never used directly: it is
recommended to use the sort \textit{int} which is no longer a builtin
data type.


\subsection{Input/Output facilities}
\label{IO switches}
%%%%%%%%%%%%%%%%%%%%%%%%%%%%%%%%%%%%%%%%%%%%%%%%%%%%%%%%%%%%
Input/Output facilities are now fully  implemented in the compiler. 
Their syntax is described in the module {\bf IO.eln} of the standard library. 

As in the version V3.4, it is possible, in the generated 
program, to enter a query term in the \elan\ syntax, just like in the
interpreter. In the same way, result terms computed by the generated
program are now pretty-printed in the \elan\ syntax. Therefore, it is
no more necessary to give a fixed ground term in the \verb|.lgi| file, 
and \verb|.lgi| files involving the \verb|query| keyword can now be
compiled. In the current version, the query is parsed dynamically in 
the generated program by reusing a subpart of the parser integrated to
the interpreter. 

The generated program can be executed with several switches used to
specify the different possible formats of input (query) and output
(result) terms:

\paragraph{Input switches}

\begin{itemize}

\item \verb|-REFInput|: the query term must be in the \REF\ 
  format (see Section~\ref{ref_format}). 

\item \verb|-noInput|: the query ground term is supposed to 
  be located in the \verb|.lgi| file, where it replaces the
  \verb|query| keyword. 


\end{itemize}



\paragraph{Output switches}

\begin{itemize}

\item \verb|-REFOutput|: result terms will be given in the
  \REF\ format (see Section~\ref{ref_format}). 

\item \verb|-noOutput|: results are computed but not
  written out. The interest of this switch is rather limited, but it
  can be useful for people who only want the statistics of the
  computation, like the number of rewrite steps per seconds. 

\item \verb|-internalOutput|: result terms are given by using a prefix 
  notation. This is the original syntax used in the first version of
  the compiler. 

\end{itemize}

The generated program, called by default \textit{a.out}, can 
be executed with some other switches not necessarily related to
Input/Ouput, like \verb|a.out -quiet| or \verb|a.out -debug|.


At this point, it is still possible to change the content of the
\verb|start with| instruction declared in the \verb|.lgi| file. The switch
\verb|-strategy|~\nt{strategy name}~{\bf :}~\nt{sort name} will apply the
strategy \nt{strategy name} on a query term of sort \nt{sort
  name}. 


Similarly, the switch \verb|-sort|~\nt{sort name} will apply 
the empty strategy on a query term \nt{sort name}. Note that
\verb|-strategy| and \verb|-sort| are mutually exclusive. 

\Attention{
The list of available sorts (resp. strategy names) can be found in the
\texttt{.ref} file.
} 


%%%%%%%%%%%%%%%%%%%%%%%%%%%%%%%%%%%%%%%%%%%%%%%%%%%%%%%%%%%%
\section{The \REF\ format}
%%%%%%%%%%%%%%%%%%%%%%%%%%%%%%%%%%%%%%%%%%%%%%%%%%%%%%%%%%%%
\label{ref_format}


The Reduced \elan\ Format, \REF\ for short, is a term-like
representation of \elan\ programs introduced both for the
interconnection of software components involved in the \elan\ system,
and for combining \elan\ programs. This format is also of greatest
interest for the implementation of reflection techniques in \elan. 
A detailed description of its applications can be found
in~\cite{BorovanskyKKMR-WRLA98}. 

The \REF\ format has been initially developed to reuse the powerful
Earley's parser integrated in the interpreter. By loading with the
interpreter an \elan\ program possibly in mixfix notation, we can get a
\REF\ program which can be viewed as the external representation of
the interpreter working memory. \REF\ is easier to parse than \elan\
since the \REF\ grammar is simply LALR.  
The \elan\ interpreter generates a \REF\ file
\textit{prog.ref} thanks to the option \texttt{--export}:
\begin{verbatim}
elan --export prog.ref prog.lgi [prog.spc]
\end{verbatim}
The \REF\ program contained in \textit{prog.ref} can be directly
executed by the interpreter as follows:
\begin{verbatim}
elan --import prog.ref
\end{verbatim}
Then, the behavior of the interpreter is exactly the one expected
with the command:
\begin{verbatim}
elan prog.lgi [prog.spc]
\end{verbatim}
but without the time-consuming preprocessing and parsing phases.

The core of the \elan\ compiler, called REM for ``Reduced \elan\ Machine'' and
implemented in \texttt{Java}, only knows the \REF\ syntax. This
explains why we need a script \texttt{elanc} to first call the
interpreter in order to obtain a \REF\ file, and second to execute REM
with this \REF\ file as input. 

Then, a compiled \texttt{C} program, 
say \texttt{a.out}, is generated by REM. This program can be executed
with different options like \texttt{a.out -REFInput}, \texttt{a.out -REFOutput}
or even better \texttt{a.out -REFInput -REFOutput}:
\begin{itemize}
\item The option \texttt{-REFInput} reads a new query in \REF\ syntax to replace
the ground query occurring initially in \textit{prog.lgi}.
\item The option \texttt{-REFOutput} writes all result terms computed by
\textit{a.out} in \REF\ syntax, instead of the \elan\ syntax used by
default. 
\end{itemize}

So, we still have to convert an \elan\ 
query to a \REF\ query, and conversely to retrieve the \elan\ syntax
of result terms written in \REF\ syntax. 

For this purpose, one can use another
tool to read a query and to pretty-print result terms of the \textit{a.out}
program, both in \elan\ syntax, similarly to the interpreter.
This tool is called \texttt{prettyIO} and may be used as follows:
\begin{verbatim}
prettyIO a.out prog.ref
\end{verbatim}
Note that \textit{prog.ref} is needed to find the \elan\ syntax of
functions implemented by \textit{a.out}, and so to encode a query
term provided by the user in \elan\ syntax, and to decode terms
written in \REF\ syntax by \texttt{a.out -REFInput -REFOutput}. 
The executable \texttt{prettyIO} is implemented as a script calling
some basic tools:  

\begin{itemize}

\item \texttt{query2ref prog.ref} reads repeatedly (from
\texttt{stdin}) a sort, and parses a term of this sort written in 
\elan\ syntax. The result (written to \texttt{stdout}) 
is the corresponding term in \REF\ syntax.
By using the option \texttt{--query} as follows: 
\begin{verbatim}
query2ref prog.ref --query
\end{verbatim}
we only parse terms of the sort of the query, as indicated in
\textit{prog.lgi} (this information also occurs in \textit{prog.ref}).

\item \texttt{ref2result prog.ref} reads repeatedly (from \texttt{stdin})
terms encoded  in \REF\ syntax and writes (to \texttt{stdout}) 
the related terms in \elan\ syntax, by using the signature 
occurring in \textit{prog.ref}.

\end{itemize}

The script \texttt{prettyIO} is located in the directory \texttt{bin/}
like \texttt{elanc}, and the compiled \texttt{C/C++} programs
\texttt{query2ref} and \texttt{ref2result} are located in the
subdirectory of \texttt{bin/} related to the current system
architecture, like the interpreter \texttt{elan}.


\Attention{
The user should be aware that \texttt{prettyIO} is much 
less interesting in the current version, since it is now possible to
parse a query term and to pretty-print all results by simply using the
generated program \textit{a.out} itself. 
}


%%% Local Variables: 
%%% mode: latex
%%% TeX-master: t
%%% TeX-master: t
%%% End: 
